\usepackage[a4paper,inner=26mm,outer=24mm,top=23mm,bottom=25mm,marginparwidth=15mm,marginparsep=6mm]{geometry}
\usepackage{fontspec}
\defaultfontfeatures{Ligatures=TeX}
\setmainfont{TeX Gyre Pagella}
\setsansfont{TeX Gyre Heros}
\setmonofont{TeX Gyre Cursor}
\usepackage{xeCJK}
\setCJKmainfont{SimSun}[AutoFakeBold=2.2]
\setCJKsansfont{Microsoft YaHei}[AutoFakeBold=2.2]
\setCJKmonofont{Microsoft YaHei}
\xeCJKsetup{PunctStyle=kaiming}
\usepackage{microtype}
\usepackage{graphicx}
\usepackage{mathtools}
\usepackage{amsthm}
\usepackage{unicode-math}
\setmathfont{Latin Modern Math}
\usepackage{bookmark}
\usepackage{hyperref}
\hypersetup{hidelinks,unicode=true}

\setcounter{secnumdepth}{3}
\setcounter{tocdepth}{2}
\numberwithin{equation}{section}
\setlength{\parindent}{2em}
\setlength{\parskip}{0pt}
\linespread{1.20}
\renewcommand{\contentsname}{目录}
\renewcommand{\proofname}{证明}

\newcommand{\fgaunitid}{unbound}
\newcommand{\src}[1]{%
  \phantomsection
  \hypertarget{src:\fgaunitid:#1}{}%
  \hypertarget{fga.\fgaunitid.p#1}{}%
  \ifcsname fgasrcseen#1\endcsname\else
    \expandafter\gdef\csname fgasrcseen#1\endcsname{1}%
    \hypertarget{src:#1}{}%
  \fi
  \marginpar{\raggedleft\footnotesize\sffamily #1}%
}
\newcommand{\sid}[1]{%
  \gdef\fgaunitid{#1}%
  \phantomsection\hypertarget{fga:#1}{}\hypertarget{fga.#1}{}%
}
\newcommand{\Osh}{\underline{O}}
\newcommand{\Zsh}{\underline{\mathbf Z}}
\newcommand{\Ccat}{\mathcal C}
\newcommand{\Homsh}{\underline{\operatorname{Hom}}}
\newcommand{\Extsh}{\underline{\operatorname{Ext}}}
\newcommand{\hsh}{\underline{h}}
\newcommand{\fquote}[1]{“#1”}
\newcommand{\zhterm}[1]{\textbf{#1}}
\newcommand{\seminarzh}[1]{%
  \begin{flushleft}\sffamily 布尔巴基讨论班\\#1\end{flushleft}%
}
\newcommand{\paperhead}[2]{%
  \begin{center}
    {\Large\bfseries #1\par}
    \vspace{0.6\baselineskip}
    #2
  \end{center}
  \vspace{0.8\baselineskip}
}

\theoremstyle{plain}
\newtheorem{proposition}{命题}[section]
\newtheorem{theoreme}[proposition]{定理}
\newtheorem{theorem}[proposition]{定理}
\newtheorem{corollaire}[proposition]{推论}
\newtheorem{corollary}[proposition]{推论}
\newtheorem{lemme}[proposition]{引理}
\newtheorem{lemma}[proposition]{引理}

\theoremstyle{definition}
\newtheorem{definition}[proposition]{定义}
\newtheorem{definitionen}[proposition]{定义}

\theoremstyle{remark}
\newtheorem*{remarque}{注}
\newtheorem*{remark}{注}

\newcommand{\unit}[1]{%
  \chapter*{#1}%
  \addcontentsline{toc}{chapter}{#1}%
  \markboth{#1}{#1}%
}
\newcommand{\fsection}[2]{%
  \section*{#1. #2}%
  \addcontentsline{toc}{section}{#1. #2}%
}
\newcommand{\result}[2]{%
  \par\medskip\noindent\textbf{#1%
  \if\relax\detokenize{#2}\relax\else\space#2\fi。——}\enspace
}
